% !TeX root = ../main.tex
% 来源：群论.md 原第二章及群论ppt第三章。沿用原稿术语、记号，修正明显笔误。
% 原稿未完成的平移群 n 维表示及 D_4 例子在正文中标作补充。
\chapter{群的线性表示理论}

本章沿用 $G\simeq D\left(G\right)$ 表示从群 $G$ 到其表示矩阵群的满同态，
$E$ 表示单位矩阵，$I$ 表示同态核。除明确讨论实表示的地方外，
表示空间均取有限维复向量空间；这一假设尤其用于舒尔定理和不可约表示的分类。

\section{线性空间和线性算符预备}
\subsection{线性空间、子空间与换基}
设 $L$ 是复数域上的 $m$ 维线性空间，$\left\{e_1,\ldots,e_m\right\}$ 是一组基，每个矢量唯一写成 $a=a_\mu e_\mu$。若
\begin{equation}
L=L_1+L_2,\qquad L_1\cap L_2=\left\{0\right\},
\end{equation}
则 $L=L_1\oplus L_2$。若线性算符 $A$ 满足 $AL_1\subseteq L_1$，则 $L_1$ 是 $A$ 的不变子空间。对群表示，要求它对所有 $D\left(g\right)$ 都不变。

设新旧基满足 $e'_\nu=e_\mu S_{\mu\nu}$，则同一矢量的分量和同一算符的矩阵满足
\begin{align}
a'&=S^{-1}a,\\
D'\left(A\right)&=S^{-1}D\left(A\right)S.
\end{align}
所以迹、行列式和特征多项式在换基下不变，等价表示正是同一组线性算符在不同基下的矩阵形式。


\section{群的线性表示}
\subsection{群的线性表示的定义}
\begin{definition}[线性表示]
设 $G=\left\{g_i\right\}$ 为抽象群，给每个群元指定一个可逆的 $m$ 阶矩阵 $D\left(g_i\right)$，满足
\begin{equation}
D\left(g_i g_j\right)=D\left(g_i\right)D\left(g_j\right).
\end{equation}
则 $D\left(G\right)=\left\{D\left(g_i\right)\mid g_i\in G\right\}$ 构成矩阵群，且 $G\simeq D\left(G\right)$。
称 $D\left(G\right)$ 为 $G$ 的一个 $m$ 维线性表示，$D\left(g_i\right)$ 为相应的表示矩阵。
\end{definition}
若对应 $g_i\mapsto D\left(g_i\right)$ 为一一对应，即由此对应给出 $G\cong D\left(G\right)$，
则沿用原稿称其为\textbf{真实表示}；否则称为\textbf{非真实表示}。
对非真实表示，同态核 $I=\left\{g\in G\mid D\left(g\right)=E\right\}$ 不仅含恒元，且
\begin{equation}
G/I\cong D\left(G\right).
\end{equation}
因此非真实表示反映的是这个商群的乘法结构。

矩阵群通过 $g\mapsto g$ 得到\textbf{自身表示}。
将全部群元映为单位矩阵，即 $D\left(G\right)=\left\{E\right\}$，得到\textbf{恒等表示}；
一维恒等表示又称平庸表示。自身表示一般不能称为平庸表示。

群元 $g_i$ 在表示中的\textbf{特征标}定义为
\begin{equation}
\chi\left(g_i\right)=\mathrm{Tr}\,D\left(g_i\right).
\end{equation}
若所有表示矩阵都满足 $D^\dagger\left(g_i\right)D\left(g_i\right)=E$，称为\textbf{幺正表示}。
若所有表示矩阵为实矩阵并满足 $D^{\mathrm{T}}\left(g_i\right)D\left(g_i\right)=E$，称为\textbf{正交表示}。

\subsection{线性表示的性质}
恒元的表示矩阵是单位矩阵。由 $D\left(e\right)D\left(g_i\right)=D\left(g_i\right)$，利用 $D\left(g_i\right)$ 可逆可得
$D\left(e\right)=E$。此外，
\begin{equation}
D\left(g_i\right)D\left(g_i^{-1}\right)=D\left(g_ig_i^{-1}\right)=D\left(e\right)=E,
\end{equation}
故 $D\left(g_i^{-1}\right)=D^{-1}\left(g_i\right)$。表示不仅保持乘法，也保持恒元和取逆。

\subsection{等价表示}
\begin{definition}[等价表示]
若存在一个与 $g_i$ 无关的非奇异矩阵 $X$，使
\begin{equation}
\bar D\left(g_i\right)=X^{-1}D\left(g_i\right)X\qquad\left(\forall g_i\in G\right),
\end{equation}
则称 $D\left(G\right)$ 与 $\bar D\left(G\right)$ 等价。
\end{definition}
这里必须对所有群元使用\textbf{同一个} $X$。等价表示是同一个线性作用在不同基下的矩阵描述。

先验证 $\bar D\left(G\right)$ 是表示：
\begin{align}
\bar D\left(g_i\right)\bar D\left(g_j\right)
&=X^{-1}D\left(g_i\right)XX^{-1}D\left(g_j\right)X\\
&=X^{-1}D\left(g_ig_j\right)X=\bar D\left(g_ig_j\right).
\end{align}
再考虑矩阵群间的对应 $A\mapsto X^{-1}AX$。
它的逆映射为 $B\mapsto XBX^{-1}$，故是双射；
且 $X^{-1}ABX=\left(X^{-1}AX\right)\left(X^{-1}BX\right)$，故保持乘法。
因此 $D\left(G\right)\cong\bar D\left(G\right)$。但仅有像矩阵群抽象同构，并不足以推出两个表示等价。

\subsection{可约表示和不可约表示}
\begin{definition}[可约表示]
设 $D\left(G\right)$ 为 $m$ 维表示。若存在非奇异矩阵 $X$，使所有表示矩阵同时化为
\begin{equation}
\bar D\left(g_i\right)=X^{-1}D\left(g_i\right)X
=\begin{bmatrix}D_1\left(g_i\right)&M\left(g_i\right)\\\bm{0}&D_2\left(g_i\right)\end{bmatrix},
\end{equation}
其中 $D_1$ 为 $k\times k$ 矩阵，$D_2$ 为 $l\times l$ 矩阵，
$M$ 为 $k\times l$ 矩阵，$k,l>0$ 且 $k+l=m$，则称该表示可约；否则称为不可约表示。
\end{definition}
前 $k$ 个基矢张成一个非零真不变子空间，即任意 $D\left(g\right)$ 将这个子空间映回自身。
反过来，先选取不变子空间的一组基，再补齐全空间的基，就得到上述分块上三角形式。
\textbf{可约不一定意味着能够化成分块对角形式}；后者还要求存在不变的补空间。

两个对角块分别给出 $G$ 的表示。事实上，
\begin{align}
\bar D\left(g_i\right)\bar D\left(g_j\right)
&=\begin{bmatrix}
D_1\left(g_i\right)D_1\left(g_j\right)&D_1\left(g_i\right)M\left(g_j\right)+M\left(g_i\right)D_2\left(g_j\right)\\
\bm{0}&D_2\left(g_i\right)D_2\left(g_j\right)
\end{bmatrix}\\
&=\bar D\left(g_ig_j\right).
\end{align}
比较对角块得到 $D_\alpha\left(g_ig_j\right)=D_\alpha\left(g_i\right)D_\alpha\left(g_j\right)$，$\alpha=1,2$。
恒元对应单位矩阵，逆元也按块满足取逆关系，所以 $G\simeq D_\alpha\left(G\right)$。
其中 $D_1$ 是不变子空间上的表示，$D_2$ 是相应商空间上的表示。

\subsection{可约表示的约化}
如果可以反复分出不变补空间，直到所有对角块不可约，则称表示\textbf{完全可约}。
记不等价不可约表示为 $\tilde D_k\left(G\right)$，其重数为 $n_k$，则
\begin{equation}
X^{-1}D\left(g\right)X=\bigoplus_{k=1}^{r}n_k\tilde D_k\left(g\right),\qquad
A\oplus B=\begin{bmatrix}A&0\\0&B\end{bmatrix}.
\end{equation}
这里 $n_k\tilde D_k$ 表示 $n_k$ 个相同矩阵块的直和，不是矩阵数乘；
例如对角线上会连续出现 $n_1$ 个 $\tilde D_1\left(g\right)$，再出现 $n_2$ 个 $\tilde D_2\left(g\right)$。
有限群的复表示总能完全约化，证明见第 2.3 节；一般无限群不能直接使用这一结论。

\subsection{\texorpdfstring{$\mathrm{SO}\left(2\right)$}{SO(2)} 群的表示}
一维表示为
\begin{equation}
D_1\left(g\left(\theta\right)\right)\equiv D_1\left(\theta\right)=\mathrm{e}^{\mathrm{i}\theta}.
\end{equation}
一维空间不存在非零真子空间，因此它不可约。
二维表示为
\begin{equation}
D_2\left(\theta\right)=\begin{bmatrix}\cos\theta&-\sin\theta\\\sin\theta&\cos\theta\end{bmatrix}.
\end{equation}
在复数域上，取基矢 $\left(1,-\mathrm{i}\right)^{\mathrm{T}}$、$\left(1,\mathrm{i}\right)^{\mathrm{T}}$，即取
\begin{equation}
X=\begin{bmatrix}1&1\\-\mathrm{i}&\mathrm{i}\end{bmatrix},\qquad
X^{-1}D_2\left(\theta\right)X=\operatorname{diag}\left(\mathrm{e}^{\mathrm{i}\theta},\mathrm{e}^{-\mathrm{i}\theta}\right),
\end{equation}
可见它完全可约。若只允许实不变子空间，则该二维实表示不可约；可约性依赖所选数域。

沿用 $\vec{\mathrm{i}},\vec{\mathrm{j}},\vec{\mathrm{k}}$ 表示三个坐标轴方向，三个三维表示为
\begin{align}
R_{\vec{\mathrm{k}}}\left(\theta\right)&=\begin{bmatrix}\cos\theta&-\sin\theta&0\\\sin\theta&\cos\theta&0\\0&0&1\end{bmatrix},\\
R_{\vec{\mathrm{j}}}\left(\theta\right)&=\begin{bmatrix}\cos\theta&0&\sin\theta\\0&1&0\\-\sin\theta&0&\cos\theta\end{bmatrix},\\
R_{\vec{\mathrm{i}}}\left(\theta\right)&=\begin{bmatrix}1&0&0\\0&\cos\theta&-\sin\theta\\0&\sin\theta&\cos\theta\end{bmatrix}.
\end{align}
每个表示都固定转轴方向，并保持垂直于转轴的平面，因而都是可约表示。

\subsection{平移群的表示}
一维平移群满足 $T\left(a\right)T\left(b\right)=T\left(a+b\right)$。沿用原稿的表示编号，一维表示有
\begin{equation}
D_1\left(a\right)=1,\qquad D_2\left(a\right)=\mathrm{e}^a,\qquad D_3\left(a\right)=\mathrm{e}^{ca},
\end{equation}
其中 $c$ 是固定复常数。它们都满足 $D_j\left(a+b\right)=D_j\left(a\right)D_j\left(b\right)$。
二维表示为
\begin{equation}
D_5\left(a\right)=\begin{bmatrix}1&a\\0&1\end{bmatrix},\qquad
D_5\left(a\right)D_5\left(b\right)=D_5\left(a+b\right).
\end{equation}
它有不变子空间 $\operatorname{span}\left\{\left(1,0\right)^{\mathrm{T}}\right\}$，但没有不变的补空间。
若补空间由 $\left(u,v\right)^{\mathrm{T}}$ 张成，其中 $v\ne0$，则
$\left(u+av,v\right)^{\mathrm{T}}$ 对所有 $a$ 都必须与 $\left(u,v\right)^{\mathrm{T}}$ 成比例。
第二分量迫使比例为 $1$，第一分量却要求 $av=0$ 对所有 $a$ 成立，矛盾。
因此该表示可约而不完全可约。

\begin{remark}[补充原稿未完成的 $n$ 维例子]
原稿的 $D_6\left(a\right)$ 矩阵为空。可取 $N$ 为上副对角线全为 $1$、其他元素为 $0$ 的
$n\times n$ 矩阵，定义
\begin{equation}
D_6\left(a\right)=\mathrm{e}^{aN}=\sum_{j=0}^{n-1}\frac{a^jN^j}{j!},\qquad
\left[D_6\left(a\right)\right]_{pq}=\begin{cases}a^{q-p}/\left(q-p\right)!,&q\ge p,\\0,&q<p.\end{cases}
\end{equation}
因为 $N^n=0$，指数级数在此截断；$aN$ 与 $bN$ 对易，故
$D_6\left(a\right)D_6\left(b\right)=D_6\left(a+b\right)$。$n=2$ 时恰好回到 $D_5\left(a\right)$。
\end{remark}

\subsection{有限循环群的表示}
设 $\mathrm{C}_n=\left\{a,a^2,\ldots,a^n=e\right\}$。一维表示由生成元的像决定：
\begin{equation}
D_k\left(a\right)=\mathrm{e}^{\mathrm{i} k2\pi/n},\qquad
D_k\left(a^j\right)=\mathrm{e}^{\mathrm{i} kj2\pi/n}.
\end{equation}
$D_0\left(a\right)=1$ 为恒等表示，$D_1$ 为真实表示；一般 $D_k$ 为真实表示的条件是
$\gcd\left(k,n\right)=1$。下标按模 $n$ 理解，$D_n=D_0$。

原稿的 $n$ 维表示来自坐标循环。记 $S=\bar D_1\left(a\right)$，规定
\begin{equation}
S\left(P_1,P_2,\ldots,P_n\right)^{\mathrm{T}}
=\left(P_n,P_1,\ldots,P_{n-1}\right)^{\mathrm{T}},
\end{equation}
即 $S_{1n}=1$、$S_{j+1,j}=1$（$1\leqslant j<n$），其余矩阵元为零。
于是 $S^n=E$，并可定义
\begin{align}
\bar D_1\left(a^i\right)&=S^i,\\
\bar D_2\left(a\right)&=S^2,\qquad \bar D_2\left(a^i\right)=S^{2i},\\
\bar D_k\left(a\right)&=S^k,\qquad \bar D_k\left(a^i\right)=S^{ki}.
\end{align}
对 $1\leqslant k<n$，按前 $k$ 行与后 $n-k$ 行分块，有
\begin{equation}
S^k=\begin{bmatrix}
\bm{0}_{k\times\left(n-k\right)}&E_{k\times k}\\
E_{\left(n-k\right)\times\left(n-k\right)}&\bm{0}_{\left(n-k\right)\times k}
\end{bmatrix}.
\end{equation}
$D_k\left(\mathrm{C}_n\right)$ 和 $\bar D_k\left(\mathrm{C}_n\right)$ 的像群均为
$n/\gcd\left(n,k\right)$ 阶循环群，因此同构；但 $n>1$ 时二者维数不同，故不等价。
$D_k$ 是不可约的一维表示，而 $\bar D_k$ 在复数域上可约。
例如 $S$ 的不同特征值为全部 $n$ 次单位根，相应特征矢量也使 $S^k$ 对角化。
$D_1,\ldots,D_n$ 恰好给出全部 $n$ 个不等价不可约复表示；完备性将在特征标理论中得到说明。

\section{标量函数的变换算符与表示的构造}
\subsection{标量函数与坐标变换}
标量场为每个空间点赋予一个标量值，可用坐标函数 $\psi\left(x\right)$ 描述。
标量性质意味着同一物理点的值不随坐标描述改变，即
\begin{equation}
x'=Rx,\qquad \psi'\left(x'\right)=\psi\left(x\right).
\end{equation}
这并不要求在同一个数值自变量处总有 $\psi'\left(x\right)=\psi\left(x\right)$。
可以保持坐标轴不动而主动变换场，也可以保持场不动而被动变换坐标；
两种描述使用互逆的坐标关系。本节用上述关系固定约定。

\subsection{由线性坐标变换诱导的变换算符}
设 $R$ 为可逆线性坐标变换，定义 $P_R\psi=\psi'$。由 $x=R^{-1}x'$ 得
\begin{equation}
\left(P_R\psi\right)\left(x\right)=\psi'\left(x\right)=\psi\left(R^{-1}x\right).
\end{equation}
这里先让算符作用于整个函数，再在 $x$ 处取值。特别地，
\begin{equation}
\left(P_R\psi\right)\left(Rx\right)=\psi\left(x\right),\qquad
\left(P_{R^{-1}}P_R\psi\right)\left(x\right)=\left(P_R\psi\right)\left(Rx\right)=\psi\left(x\right).
\end{equation}
交换 $R$ 与 $R^{-1}$ 也成立，因此 $P_{R^{-1}}=P_R^{-1}$。
求变换后的函数，就是将原函数的自变量替换为 $R^{-1}x$。

\begin{example}[反射对标量函数的作用]
取 $\psi\left(\vec r\right)=xy$，$R:\left(x,y,z\right)\mapsto\left(-x,y,z\right)$。
这个变换满足 $R^{-1}=R$，故
\begin{equation}
\psi'\left(\vec r\right)=\left(P_R\psi\right)\left(\vec r\right)=\psi\left(R^{-1}\vec r\right)=\left(-x\right)y=-xy.
\end{equation}
\end{example}

\subsection{坐标变换群与函数变换算符群的同构}
首先，$P_R$ 为线性算符。对常数 $a,b$，
\begin{align}
\left(P_R\left(a\psi+b\phi\right)\right)\left(x\right)
&=a\psi\left(R^{-1}x\right)+b\phi\left(R^{-1}x\right)\\
&=a\left(P_R\psi\right)\left(x\right)+b\left(P_R\phi\right)\left(x\right).
\end{align}
其次，乘法次序满足
\begin{align}
\left(P_SP_R\psi\right)\left(x\right)&=\left(P_R\psi\right)\left(S^{-1}x\right)\\
&=\psi\left(R^{-1}S^{-1}x\right)=\psi\left(\left(SR\right)^{-1}x\right)=\left(P_{SR}\psi\right)\left(x\right).
\end{align}
所以 $P_{SR}=P_SP_R$，即 $R\mapsto P_R$ 给出 $\left\{R\right\}\simeq\left\{P_R\right\}$。

若函数空间包含全部坐标函数，则此对应还是单射。
假设 $P_R=P_S$，将等式作用于每个坐标函数 $x_j$，就有
$\left(R^{-1}x\right)_j=\left(S^{-1}x\right)_j$ 对所有 $j,x$ 成立，故 $R^{-1}=S^{-1}$，即 $R=S$。
因此在这样的函数空间上 $\left\{R\right\}\cong\left\{P_R\right\}$。
只对某一个函数验证 $\psi\left(R^{-1}x\right)=\psi\left(S^{-1}x\right)$ 不足以推出 $R=S$；
限制到某个子空间后，所得到的表示也可能是非真实表示。

\subsection{利用基函数构造线性表示}
取一个 $n$ 维函数空间，对所有 $P_R$ 都不变，基函数排为行矢量
$\psi=\left(\psi_1,\ldots,\psi_n\right)$。定义矩阵元
\begin{equation}
P_R\psi_\mu=\sum_{\nu=1}^{n}\psi_\nu D_{\nu\mu}\left(R\right),\qquad
P_R\psi=\psi D\left(R\right).
\end{equation}
矩阵 $D\left(R\right)$ 的第 $\mu$ 列就是 $P_R\psi_\mu$ 在所选基下的展开系数；
例如其第二行第二列记为 $D_{22}\left(R\right)$。
因为系数是常数，
\begin{equation}
P_SP_R\psi=P_S\left(\psi D\left(R\right)\right)=\left(P_S\psi\right)D\left(R\right)=\psi D\left(S\right)D\left(R\right).
\end{equation}
另一方面 $P_{SR}\psi=\psi D\left(SR\right)$。利用基函数线性无关，得到
\begin{equation}
D\left(SR\right)=D\left(S\right)D\left(R\right).
\end{equation}
因此 $\left\{P_R\right\}\simeq\left\{D\left(R\right)\right\}$，也得到 $\left\{R\right\}\simeq\left\{D\left(R\right)\right\}$。
表示的维数就是基函数个数。

若某个非零真子空间不变，则表示可约。若它存在不变补空间，则能分成两个对角块；
进一步完全分解为不可约块时，表示完全可约。
判断不完全可约时，必须确认\textbf{不存在}所需的不变补空间，
而不是仅仅发现任选的某个补空间不变性不成立。

\begin{example}[补充：$\mathrm{D}_4$ 在三维空间中的表示]
原稿只指定表示空间为 $\mathbb{R}^3$。这里选取正方形位于 $xy$ 平面的平面对称群，
令所有变换保持 $z$ 不变。取生成元 $r$ 为绕 $z$ 轴转动 $\pi/2$，$s$ 为关于 $xz$ 平面反射：
\begin{equation}
r=\begin{bmatrix}0&-1&0\\1&0&0\\0&0&1\end{bmatrix},\qquad
s=\begin{bmatrix}1&0&0\\0&-1&0\\0&0&1\end{bmatrix}.
\end{equation}
它们满足 $r^4=s^2=e$、$srs=r^{-1}$，生成八个元素 $r^j,sr^j$（$0\leqslant j<4$）。
在基函数 $\left(x,y,z\right)$ 下，$P_r\left(x,y,z\right)=\left(y,-x,z\right)$、$P_s\left(x,y,z\right)=\left(x,-y,z\right)$，
所以 $D\left(r\right)=r$、$D\left(s\right)=s$。这相当于将上述三维坐标空间与线性函数空间对应起来。
空间分解为 $\operatorname{span}\left\{x,y\right\}\oplus\operatorname{span}\left\{z\right\}$；
后一项为一维恒等表示。前一项在复数域上也不可约：$r$ 的两个特征直线被 $s$ 互换，
没有同时对两者不变的直线。
\end{example}

\section{有限群表示理论的基本定理}
本节设 $G$ 为有限群，阶为 $n_G$，表示均为有限维复表示。

\begin{theorem}[定理 1：幺正化]
有限群的任意线性表示均等价于一个幺正表示。
\end{theorem}
\begin{proof}
要使 $\bar D\left(g\right)=X^{-1}D\left(g\right)X$ 幺正，等价于要求
\begin{equation}
D^\dagger\left(g\right)\left(XX^\dagger\right)^{-1}D\left(g\right)=\left(XX^\dagger\right)^{-1}.
\end{equation}
定义 $H=\sum_{f\in G}D^\dagger\left(f\right)D\left(f\right)$。它是 Hermitian 矩阵，且任意非零列向量 $v$ 满足
\begin{equation}
v^\dagger Hv=\sum_{f\in G}\|D\left(f\right)v\|^2\geqslant\|D\left(e\right)v\|^2=\|v\|^2>0,
\end{equation}
所以 $H$ 正定。由重排定理，
\begin{align}
D^\dagger\left(g\right)HD\left(g\right)
&=\sum_{f\in G}\left[D\left(f\right)D\left(g\right)\right]^\dagger\left[D\left(f\right)D\left(g\right)\right]\\
&=\sum_{f\in G}D^\dagger\left(fg\right)D\left(fg\right)=H.
\end{align}
用幺正矩阵 $U$ 对角化 $H$：
\begin{equation}
\Gamma=U^{-1}HU=\operatorname{diag}\left(\Gamma_{11},\ldots,\Gamma_{mm}\right),\qquad \Gamma_{jj}>0.
\end{equation}
令 $\Gamma'=\operatorname{diag}\left(1/\sqrt{\Gamma_{11}},\ldots,1/\sqrt{\Gamma_{mm}}\right)$，
取 $X=U\Gamma'U^{-1}$。由于 $U^{-1}=U^\dagger$，有 $X^\dagger=X$，从而
\begin{equation}
\left(XX^\dagger\right)^{-1}=U\left(\Gamma'\right)^{-2}U^{-1}=U\Gamma U^{-1}=H.
\end{equation}
因此该 $X$ 满足所需条件，完成幺正化。
\end{proof}

\begin{proposition}[完全可约性]
有限群的复表示完全可约。
\end{proposition}
\begin{proof}
先将表示幺正化。设 $W$ 为不变子空间，$v\in W^\perp$、$w\in W$，则
\begin{equation}
\langle w,D\left(g\right)v\rangle=\langle D\left(g\right)^{-1}w,v\rangle=0,
\end{equation}
因为 $D\left(g\right)^{-1}w\in W$。故 $W^\perp$ 也不变，表示可分为两个对角块。
对各块重复此过程；维数有限保证最终终止于不可约块。
\end{proof}

\begin{theorem}[定理 2：等价幺正表示]
若 $D\left(G\right)$、$\bar D\left(G\right)$ 为等价的幺正表示，则存在幺正矩阵 $U$，使
$\bar D\left(g\right)=U^{-1}D\left(g\right)U$ 对所有 $g\in G$ 成立。
\end{theorem}
\begin{proof}
设 $D\left(g\right)X=X\bar D\left(g\right)$。对 $g^{-1}$ 的等式取 Hermitian 共轭，得到
$X^\dagger D\left(g\right)=\bar D\left(g\right)X^\dagger$。
因此 $X^\dagger X$ 与所有 $\bar D\left(g\right)$ 对易，其正定平方根
$A=\left(X^\dagger X\right)^{1/2}$ 也与它们对易；后者可由正定矩阵的谱分解得到。
令 $U=XA^{-1}$，则 $U^\dagger U=A^{-1}X^\dagger XA^{-1}=E$，且
$D\left(g\right)U=X\bar D\left(g\right)A^{-1}=U\bar D\left(g\right)$，所以得到所需幺正相似变换。
\end{proof}

\subsection{定理 3：舒尔定理}
\begin{theorem}[舒尔定理一]
设 $D\left(G\right)$ 为不可约复表示。若 $D\left(g\right)X=XD\left(g\right)$ 对任意 $g\in G$ 成立，
则 $X=cE$，其中 $c$ 为复常数。
\end{theorem}
\begin{proof}
复矩阵 $X$ 至少有一个特征值 $c$。对应的非零特征子空间 $W$ 满足：
若 $Xv=cv$，则 $X D\left(g\right)v=D\left(g\right)Xv=cD\left(g\right)v$，故 $W$ 不变。
由不可约性，$W$ 必为全空间，因而 $X=cE$。
\end{proof}

\begin{theorem}[舒尔定理二]
设 $D^{\left(1\right)}\left(G\right)$、$D^{\left(2\right)}\left(G\right)$ 为不等价不可约表示，维数为 $m_1,m_2$。
若 $m_1\times m_2$ 矩阵 $X$ 满足
\begin{equation}
D^{\left(1\right)}\left(g\right)X=XD^{\left(2\right)}\left(g\right)\qquad\left(\forall g\in G\right),
\end{equation}
则 $X=\bm{0}$。
\end{theorem}
\begin{proof}
上式说明 $X$ 的核在第二个表示下不变，像在第一个表示下不变。
若 $X\ne0$，不可约性迫使其核为零、像为全空间，故 $X$ 为可逆映射，
进而给出两表示等价，与假设矛盾。
\end{proof}

\subsection{定理 4：正交性定理}
\begin{theorem}[表示矩阵元的正交性]
设 $D^{\left(u\right)}$ 为两两不等价的不可约幺正表示，维数为 $n_u$，则
\begin{equation}
\sum_{\alpha=1}^{n_G}D^{\left(u\right)*}_{\gamma\eta}\left(g_\alpha\right)
D^{\left(v\right)}_{\lambda\rho}\left(g_\alpha\right)
=\frac{n_G}{n_u}\delta^{uv}\delta_{\gamma\lambda}\delta_{\eta\rho}.
\end{equation}
\end{theorem}
\begin{proof}
取任意 $n_v\times n_u$ 矩阵 $A$，定义
\begin{equation}
B=\sum_{g\in G}D^{\left(v\right)}\left(g\right)A D^{\left(u\right)}\left(g\right)^{-1}.
\end{equation}
由重排定理，$D^{\left(v\right)}\left(h\right)B D^{\left(u\right)}\left(h\right)^{-1}=B$。
舒尔定理给出：$u\ne v$ 时 $B=0$；$u=v$ 时 $B=cE$。
后一种情况下取迹得到 $n_uc=n_G\mathrm{Tr}\,A$，所以
$B=\left(n_G/n_u\right)\left(\mathrm{Tr}\,A\right)E$。
令 $A$ 仅在第 $\rho$ 行、第 $\eta$ 列取值 $1$，取 $B$ 的第 $\lambda$ 行、第 $\gamma$ 列，
再利用 $D^{\left(u\right)}\left(g\right)^{-1}=D^{\left(u\right)}\left(g\right)^\dagger$，即得所需公式。
\end{proof}

\subsection{定理 5：不可约表示的个数}
\begin{theorem}
有限群不等价不可约复表示的个数，等于群的共轭类的个数。
\end{theorem}
本定理的证明将在第 2.5 节利用特征标的正交性与正则表示给出，
以便先定义所需对象。记这一共同数目为 $r$。

\section{群的正则表示}
\subsection{群空间}
设有限群 $G=\left\{g_1,\ldots,g_n\right\}$，其中 $n=n_G$。
将群元视为\textbf{形式上的线性无关基矢}，它们张成的复线性空间称为群空间。
记基矢行向量 $v=\left(g_1,\ldots,g_n\right)$，则一般矢量为
\begin{equation}
f=x_1g_1+\cdots+x_ng_n=vx,\qquad x=\left(x_1,\ldots,x_n\right)^{\mathrm{T}}.
\end{equation}
群元即使原本为矩阵，这里的线性组合也按形式基矢理解，不采用原矩阵之间可能存在的线性关系。

\subsection{群的正则表示}
让群元通过\textbf{左乘}作用于群空间：
\begin{equation}
g_\alpha f=\sum_{j=1}^{n}x_j\left(g_\alpha g_j\right).
\end{equation}
将结果重新在 $v$ 中展开，定义 $g_\alpha v=vD\left(g_\alpha\right)$。矩阵元为
\begin{equation}
D\left(g_\alpha\right)_{ij}=\begin{cases}1,&g_\alpha g_j=g_i,\\0,&\text{其他情形}.
\end{cases}
\end{equation}
左乘的结合律保证 $D\left(g_\alpha g_\beta\right)=D\left(g_\alpha\right)D\left(g_\beta\right)$，这就是左正则表示。
重排定理保证每行、每列恰有一个 $1$，所以它由置换矩阵组成。
对于 $f'=g_\alpha f$，可以写成
\begin{equation}
f'=vD\left(g_\alpha\right)x=v'x=vx',\qquad
v'=vD\left(g_\alpha\right),\quad x'=D\left(g_\alpha\right)x.
\end{equation}
第一种写法让各基矢随左乘变换而保持系数不变；第二种在固定基中重新展开，系数变为 $x'$。

\begin{example}[四阶群的正则表示]
取原稿的乘法表，行元在左、列元在右：
\begin{center}
\begin{tabular}{c|cccc}
\toprule
 & $g_1$ & $g_2$ & $g_3$ & $g_4$\\\midrule
$g_1$ & $g_1$ & $g_2$ & $g_3$ & $g_4$\\
$g_2$ & $g_2$ & $g_1$ & $g_4$ & $g_3$\\
$g_3$ & $g_3$ & $g_4$ & $g_2$ & $g_1$\\
$g_4$ & $g_4$ & $g_3$ & $g_1$ & $g_2$\\\bottomrule
\end{tabular}
\end{center}
设 $f=x_1g_1+x_2g_2+x_3g_3+x_4g_4$，则
\begin{align}
g_2f&=x_1g_2g_1+x_2g_2g_2+x_3g_2g_3+x_4g_2g_4\\
&=x_1g_2+x_2g_1+x_3g_4+x_4g_3\\
&=x_2g_1+x_1g_2+x_4g_3+x_3g_4.
\end{align}
故 $x'=\left(x_2,x_1,x_4,x_3\right)^{\mathrm{T}}=D\left(g_2\right)x$。
从基矢观点，$\left(g_2,g_1,g_4,g_3\right)=vD\left(g_2\right)$，得到同一个矩阵。
四个矩阵为
\begin{align}
D\left(g_1\right)&=\begin{bmatrix}1&0&0&0\\0&1&0&0\\0&0&1&0\\0&0&0&1\end{bmatrix},&
D\left(g_2\right)&=\begin{bmatrix}0&1&0&0\\1&0&0&0\\0&0&0&1\\0&0&1&0\end{bmatrix},\\
D\left(g_3\right)&=\begin{bmatrix}0&0&0&1\\0&0&1&0\\1&0&0&0\\0&1&0&0\end{bmatrix},&
D\left(g_4\right)&=\begin{bmatrix}0&0&1&0\\0&0&0&1\\0&1&0&0\\1&0&0&0\end{bmatrix}.
\end{align}
其中 $D\left(g_2\right)$ 的最后一行为 $\left(0,0,1,0\right)$，由上面的系数变换直接确定。
\end{example}

\subsection{正则表示的特点}
\begin{enumerate}
\item 若对角元 $D\left(g_\alpha\right)_{jj}=1$，则 $g_\alpha g_j=g_j$，从而 $g_\alpha=e$。
所以正则表示的特征标为
\begin{equation}
\chi\left(g_\alpha\right)=\begin{cases}n_G,&g_\alpha=e,\\0,&g_\alpha\ne e.\end{cases}
\end{equation}
\item 正则表示的维数为 $n_G$，因为群空间有 $n_G$ 个形式基矢。
\item 正则表示是真实表示：若 $D\left(g\right)=D\left(h\right)$，让它们作用于恒元基矢，得到 $g=h$。
因此 $G\cong D\left(G\right)$。
\item 任何有限群都有上述正则表示，不需要预先给出几何实现。
\end{enumerate}

\section{特征标相关理论}
本节仍讨论有限群的有限维复表示。$n_D$ 表示某个表示的维数，
$n_u$ 表示第 $u$ 个不可约表示的维数，$n_{C_\alpha}$ 表示共轭类 $C_\alpha$ 的大小。

\subsection{特征标是类的函数}
特征标为
\begin{equation}
\chi\left(g_\alpha\right)=\mathrm{Tr}\,D\left(g_\alpha\right)=\sum_{i=1}^{n_D}D_{ii}\left(g_\alpha\right).
\end{equation}
由表示保持乘法和迹的循环性质，
\begin{align}
\chi\left(f^{-1}g_\alpha f\right)
&=\mathrm{Tr}\left(D\left(f\right)^{-1}D\left(g_\alpha\right)D\left(f\right)\right)\\
&=\mathrm{Tr}\left(D\left(f\right)D\left(f\right)^{-1}D\left(g_\alpha\right)\right)=\chi\left(g_\alpha\right).
\end{align}
所以同类元素的特征标相同，特征标是类函数。

\begin{theorem}[定理 1]
等价表示具有相同的特征标。
\end{theorem}
\begin{proof}
若 $\bar D\left(g_\alpha\right)=X^{-1}D\left(g_\alpha\right)X$，则
\begin{equation}
\bar\chi\left(g_\alpha\right)=\mathrm{Tr}\left(X^{-1}D\left(g_\alpha\right)X\right)
=\mathrm{Tr}\,D\left(g_\alpha\right)=\chi\left(g_\alpha\right).
\end{equation}
\end{proof}

\subsection{定理 2：特征标正交关系}
\begin{theorem}
不等价不可约表示的特征标满足
\begin{equation}
\sum_{\alpha=1}^{n_G}\chi^{\left(u\right)*}\left(g_\alpha\right)\chi^{\left(v\right)}\left(g_\alpha\right)
=n_G\delta^{uv}.
\end{equation}
\end{theorem}
\begin{proof}
由幺正化及特征标在等价变换下不变，只需考虑幺正表示。
在矩阵元正交关系中令两个矩阵元均取对角元，得到
\begin{equation}
\sum_{\alpha=1}^{n_G}D^{\left(u\right)*}_{\sigma\sigma}\left(g_\alpha\right)
D^{\left(v\right)}_{\lambda\lambda}\left(g_\alpha\right)
=\frac{n_G}{n_u}\delta^{uv}\delta_{\sigma\lambda}\delta_{\sigma\lambda}.
\end{equation}
对 $\sigma,\lambda$ 求和，左侧成为特征标的内积。
当 $u\ne v$ 时右侧为零；当 $u=v$ 时
$\sum_{\sigma,\lambda}\delta_{\sigma\lambda}\delta_{\sigma\lambda}=n_u$，故得到结论。
\end{proof}

按群元排列的特征标表为
\begin{center}
\begin{tabular}{c|cccc}
\toprule
 & $g_1$ & $g_2$ & $\cdots$ & $g_{n_G}$\\\midrule
$D^{\left(1\right)}$ & $\chi^{\left(1\right)}\left(g_1\right)$ & $\chi^{\left(1\right)}\left(g_2\right)$ & $\cdots$ & $\chi^{\left(1\right)}\left(g_{n_G}\right)$\\
$D^{\left(2\right)}$ & $\chi^{\left(2\right)}\left(g_1\right)$ & $\chi^{\left(2\right)}\left(g_2\right)$ & $\cdots$ & $\chi^{\left(2\right)}\left(g_{n_G}\right)$\\
$\vdots$ & $\vdots$ & $\vdots$ & $\ddots$ & $\vdots$\\
$D^{\left(r\right)}$ & $\chi^{\left(r\right)}\left(g_1\right)$ & $\chi^{\left(r\right)}\left(g_2\right)$ & $\cdots$ & $\chi^{\left(r\right)}\left(g_{n_G}\right)$\\\bottomrule
\end{tabular}
\end{center}
一行对应一个不可约表示，称为特征标行矢量；一列对应一个群元，称为特征标列矢量。
任意不同行正交，每行的\textbf{模方}为 $n_G$，模长为 $\sqrt{n_G}$：
\begin{equation}
|\chi^{\left(u\right)}\left(G\right)|^2=\sum_{\alpha=1}^{n_G}|\chi^{\left(u\right)}\left(g_\alpha\right)|^2=n_G.
\end{equation}

\subsection{定理 3：唯一分解定理}
\begin{theorem}
设 $D^{\left(u\right)}$ 遍历全部不等价不可约表示，则任意表示可以分解为
\begin{equation}
X^{-1}D\left(g_\alpha\right)X=\bigoplus_{u=1}^{r}a_uD^{\left(u\right)}\left(g_\alpha\right),
\end{equation}
其中非负整数重数 $a_u$ 唯一确定，且
\begin{equation}
a_u=\frac1{n_G}\sum_{\alpha=1}^{n_G}\chi^{\left(u\right)*}\left(g_\alpha\right)\chi\left(g_\alpha\right)
=\frac1{n_G}\sum_{\alpha=1}^{n_G}\chi^*\left(g_\alpha\right)\chi^{\left(u\right)}\left(g_\alpha\right).
\end{equation}
\end{theorem}
\begin{proof}
完全可约性保证存在上述直和分解。两侧取迹得
$\chi\left(g_\alpha\right)=\sum_u a_u\chi^{\left(u\right)}\left(g_\alpha\right)$。
乘以 $\chi^{\left(v\right)*}\left(g_\alpha\right)$ 后求和，利用正交性，
\begin{align}
\sum_\alpha\chi^{\left(v\right)*}\left(g_\alpha\right)\chi\left(g_\alpha\right)
&=\sum_u a_u\sum_\alpha\chi^{\left(v\right)*}\left(g_\alpha\right)\chi^{\left(u\right)}\left(g_\alpha\right)\\
&=n_Ga_v.
\end{align}
故得到第一种写法。重数是非负整数，因而等于自身复共轭；对公式取复共轭，
即得原稿采用的第二种写法。重数由特征标确定，所以与约化过程无关。
\end{proof}

\subsubsection*{练习：可约表示特征标的模方}
利用 $a_u$ 为实数及正交性，
\begin{align}
\sum_{\beta=1}^{n_G}\chi^*\left(g_\beta\right)\chi\left(g_\beta\right)
&=\sum_{u,v}a_ua_v\sum_\beta\chi^{\left(u\right)*}\left(g_\beta\right)\chi^{\left(v\right)}\left(g_\beta\right)\\
&=n_G\sum_u a_u^2.
\end{align}
非零表示不可约，当且仅当恰有一个 $a_u=1$、其余为零。
若表示可约，则 $\sum_u a_u^2\geq2$，因而模方至少为 $2n_G$。

\subsubsection*{应用于正则表示}
正则表示仅在恒元处的特征标非零，故
\begin{equation}
a_u=\frac1{n_G}\chi^*\left(e\right)\chi^{\left(u\right)}\left(e\right)
=\frac1{n_G}n_Gn_u=n_u.
\end{equation}
因此每个不可约表示都在正则表示中出现，重数等于其维数：
\begin{equation}
X^{-1}D_{\mathrm{regular}}\left(g_\alpha\right)X
=\bigoplus_{u=1}^{r}n_uD^{\left(u\right)}\left(g_\alpha\right).
\end{equation}
在恒元处取迹得到
\begin{equation}
n_G=\sum_{u=1}^{r}n_u^2.
\end{equation}
在非恒元处则有
\begin{equation}
0=\sum_{u=1}^{r}n_u\chi^{\left(u\right)}\left(g_\alpha\right)
=\sum_{u=1}^{r}n_u\chi^{\left(u\right)*}\left(g_\alpha\right),\qquad g_\alpha\ne e.
\end{equation}
因为 $n_u=\chi^{\left(u\right)}\left(e\right)$，特征标表的恒元列与其他列正交。

\subsubsection*{补证：不可约表示的个数等于类数}
为避免先使用结论，暂记类数为 $c$、不等价不可约表示个数为 $r$。
类函数空间维数为 $c$，其内积取
$\langle F_1,F_2\rangle=n_G^{-1}\sum_g F_1\left(g\right)^*F_2\left(g\right)$。
不可约特征标为一组正交非零类函数，因此 $r\leqslant c$。

若类函数 $F$ 与所有不可约特征标正交，则
$\sum_g F\left(g\right)^*\chi^{\left(u\right)}\left(g\right)=0$。
在群空间中引入形式线性组合 $z=\sum_g F\left(g\right)^*g$，并将表示按线性延拓作用于它。
由于 $F$ 在共轭类上为常数，$z$ 与所有群元对易，所以
$D^{\left(u\right)}\left(z\right)$ 与所有 $D^{\left(u\right)}\left(g\right)$ 对易。舒尔定理给出
$D^{\left(u\right)}\left(z\right)=c_uE$，而其迹为零，故 $c_u=0$。
正则表示完全可约，故 $z$ 在正则表示中也作用为零。
让它左乘恒元基矢，得到 $ze=z=0$；由形式基矢线性无关，$F\left(g\right)=0$ 对所有 $g$ 成立。
所以不可约特征标张成全部类函数空间，$r=c$，完成第 2.3 节定理 5 的证明。

\subsection{以类为表头的特征标表}
将同类群元的列合并，约定 $D^{\left(1\right)}$ 为一维恒等表示，$C_1=\left\{e\right\}$，则
\begin{center}
\begin{tabular}{c|cccc}
\toprule
 & $C_1=\left\{e\right\}$ & $C_2$ & $\cdots$ & $C_r$\\\midrule
$D^{\left(1\right)}$ & $1$ & $1$ & $\cdots$ & $1$\\
$D^{\left(2\right)}$ & $n_2$ & $\chi^{\left(2\right)}\left(C_2\right)$ & $\cdots$ & $\chi^{\left(2\right)}\left(C_r\right)$\\
$\vdots$ & $\vdots$ & $\vdots$ & $\ddots$ & $\vdots$\\
$D^{\left(r\right)}$ & $n_r$ & $\chi^{\left(r\right)}\left(C_2\right)$ & $\cdots$ & $\chi^{\left(r\right)}\left(C_r\right)$\\\bottomrule
\end{tabular}
\end{center}
类 $C_k$ 有时在表头记作 $n_{C_k}g_k$，其中 $g_k\in C_k$ 是代表元。
这里显式写 $n_{C_k}$，避免将类大小与不可约表示维数 $n_k$ 混淆。
合并后的表为 $r\times r$ 方阵，行正交关系必须计入每类的大小：
\begin{equation}
\sum_{\alpha=1}^{r}n_{C_\alpha}\chi^{\left(u\right)*}\left(C_\alpha\right)\chi^{\left(v\right)}\left(C_\alpha\right)
=n_G\delta^{uv}.
\end{equation}

定义归一化特征标
\begin{equation}
\hat\chi^{\left(u\right)}\left(C_\alpha\right)=\sqrt{\frac{n_{C_\alpha}}{n_G}}\chi^{\left(u\right)}\left(C_\alpha\right).
\end{equation}
归一化特征标组成方阵，其行正交且模长为 $1$，因而该方阵幺正，列也正交。
因此
\begin{align}
\sum_{\alpha=1}^{r}\hat\chi^{\left(u\right)*}\left(C_\alpha\right)\hat\chi^{\left(v\right)}\left(C_\alpha\right)&=\delta^{uv},\\
\sum_{u=1}^{r}\hat\chi^{\left(u\right)*}\left(C_\alpha\right)\hat\chi^{\left(u\right)}\left(C_\beta\right)&=\delta_{\alpha\beta}.
\end{align}
去掉归一化因子，得到列正交关系
\begin{equation}
\sum_{u=1}^{r}\chi^{\left(u\right)*}\left(C_\alpha\right)\chi^{\left(u\right)}\left(C_\beta\right)
=\frac{n_G}{n_{C_\alpha}}\delta_{\alpha\beta}.
\end{equation}
若列用任意群元而非类标记，同一类的两个群元对应相同的列，不能声称它们正交。
上述 Kronecker 符号区分的是\textbf{共轭类}。

每个归一化不可约特征标行矢量和列矢量的模方均为 $1$：
\begin{equation}
\sum_\alpha|\hat\chi^{\left(u\right)}\left(C_\alpha\right)|^2=1,\qquad
\sum_u|\hat\chi^{\left(u\right)}\left(C_\alpha\right)|^2=1.
\end{equation}
对于一个待判定的非零表示 $D$，定义其行矢量
$\hat\chi\left(C_\alpha\right)=\sqrt{n_{C_\alpha}/n_G}\chi\left(C_\alpha\right)$，则
\begin{equation}
D\text{ 不可约}\quad\Longleftrightarrow\quad
\sum_{\alpha=1}^{r}|\hat\chi\left(C_\alpha\right)|^2=1.
\end{equation}
若 $D$ 可约，则该模方为 $\sum_u a_u^2\geq2$。判据应对类指标 $\alpha$ 求和，
不是对所有不可约表示指标 $u$ 求和。

\begin{example}[$\mathrm{D}_3$ 的特征标表]
$\mathrm{D}_3$ 有三个类：$\left\{e\right\}$、$\left\{d,f\right\}$、$\left\{a,b,c\right\}$，大小分别为 $1,2,3$。
由 $r=3$ 及 $n_G=6$，维数满足 $6=1^2+1^2+2^2$。
两个一维表示分别是恒等表示和将转动映为 $1$、反射映为 $-1$ 的表示。
将二维表示的两个未知特征标记为 $x,y$，有
\begin{center}
\begin{tabular}{c|ccc}
\toprule
 & $\left\{e\right\}$ & $\left\{d,f\right\}$ & $\left\{a,b,c\right\}$\\\midrule
$D^{\left(1\right)}$ & 1 & 1 & 1\\
$D^{\left(2\right)}$ & 1 & 1 & $-1$\\
$D^{\left(3\right)}$ & 2 & $x$ & $y$\\\bottomrule
\end{tabular}
\end{center}
第一列与其余两列正交，给出 $1+1+2x=0$、$1-1+2y=0$，
所以 $x=-1$、$y=0$。例如第三行的加权模方为
$1\cdot2^2+2\cdot\left(-1\right)^2+3\cdot0^2=6$，符合正交关系。
\end{example}

\begin{example}[$\mathrm{D}_4$ 的不可约表示维数]
用 $r_0,s_0$ 表示满足 $r_0^4=s_0^2=e$、$s_0r_0s_0=r_0^{-1}$ 的生成元。
五个类为
\begin{equation}
\left\{e\right\},\quad\left\{r_0^2\right\},\quad\left\{r_0,r_0^3\right\},\quad
\left\{s_0,s_0r_0^2\right\},\quad\left\{s_0r_0,s_0r_0^3\right\}.
\end{equation}
结合类数为 $5$ 与群阶为 $8$，得到
$8=1^2+1^2+1^2+1^2+2^2$，即四个一维和一个二维不等价不可约表示。
仅凭一个平方和分解，而不使用类数等信息，一般不足以确定不可约表示的维数。
\end{example}

\subsection{定理 4：特征标判定等价性}
\begin{theorem}
有限群的两个复表示等价，当且仅当它们对每个群元具有相同的特征标。
\end{theorem}
\begin{proof}
等价推出特征标相同已在定理 1 中证明。
反过来，设 $\bar\chi\left(g_\alpha\right)=\chi\left(g_\alpha\right)$。将两个表示完全约化：
\begin{align}
M^{-1}\bar D\left(g_\alpha\right)M&=\bigoplus_u\bar a_uD^{\left(u\right)}\left(g_\alpha\right),\\
N^{-1}D\left(g_\alpha\right)N&=\bigoplus_u a_uD^{\left(u\right)}\left(g_\alpha\right).
\end{align}
唯一分解定理给出
\begin{equation}
\bar a_u=\frac1{n_G}\sum_\alpha\bar\chi^*\left(g_\alpha\right)\chi^{\left(u\right)}\left(g_\alpha\right)
=\frac1{n_G}\sum_\alpha\chi^*\left(g_\alpha\right)\chi^{\left(u\right)}\left(g_\alpha\right)=a_u.
\end{equation}
按相同次序排列不可约块，则 $M^{-1}\bar D\left(g_\alpha\right)M=N^{-1}D\left(g_\alpha\right)N$。
令 $X=NM^{-1}$，得到 $\bar D\left(g_\alpha\right)=X^{-1}D\left(g_\alpha\right)X$，所以两表示等价。
\end{proof}

\subsection{特征标表的性质与应用}
\begin{enumerate}
\item 以类为表头的不可约特征标表为 $r\times r$ 方阵。
\item 归一化后任意两行正交：$\sum_\alpha\hat\chi^{\left(u\right)*}\left(C_\alpha\right)\hat\chi^{\left(v\right)}\left(C_\alpha\right)=\delta^{uv}$。
\item 归一化后任意两列正交：$\sum_u\hat\chi^{\left(u\right)*}\left(C_\alpha\right)\hat\chi^{\left(u\right)}\left(C_\beta\right)=\delta_{\alpha\beta}$。
\item 归一化后的每行、每列模方均为 $1$。
\item 恒元列给出 $n_G^{-1}\sum_u n_u^2=1$。
\item 商群的不可约表示可给出原群的不可约表示。
设 $H$ 为 $G$ 的不变子群，$D'$ 为 $G/H$ 的不可约表示，令 $D\left(g\right)=D'\left(gH\right)$。
因为 $gH$ 遍历整个商群，两表示的不变子空间相同，故 $D$ 不可约。
因此可以先找不变子群与商群，再构造原群的一部分不可约表示。
\item 一维非恒等表示与高维不可约表示逐群元相乘，仍得不可约表示。
若一维表示为 $d\left(g\right)$，则 $\bar D\left(g\right)=d\left(g\right)D\left(g\right)$ 保持乘法。
因 $d\left(g\right)\ne0$ 为标量，$\bar D$ 与 $D$ 的不变子空间相同，故不可约性不变。
该结论对一维恒等表示也成立。
\item 若 $G=H_1\otimes H_2$，则每个类是两个因子中相应类的乘积，
所以 $G$ 的类数等于 $H_1$ 与 $H_2$ 类数之积。
\end{enumerate}

\begin{example}[循环群]
沿用此处原稿的生成元记号，$\mathrm{C}_n=\left\{e,c_n,\ldots,c_n^{n-1}\right\}$。
它是 Abel 群，每个群元自成一类，因此有 $n$ 个不等价不可约表示。
由 $\sum_{u=1}^{n}n_u^2=n$ 及 $n_u\geq1$，它们全部为一维，具体为
\begin{equation}
D^{\left(k\right)}\left(c_n\right)=\mathrm{e}^{\mathrm{i}2\pi k/n},\qquad k=1,\ldots,n.
\end{equation}
这些数两两不同，故表示不等价；它们已经穷尽全部不可约表示。
特别地，$\mathrm{C}_6=\mathrm{C}_3\otimes\mathrm{C}_2$，与第一章的直积分解一致。
\end{example}

\begin{example}[二面体群]
沿用原稿的 $\mathrm{D}_n$ 表示 $2n$ 阶二面体群。
在 $\mathrm{D}_2=\left\{e,c_2=\sigma_x,c_2'=\sigma_y,c_2''=\sigma_z\right\}$ 中，
$\sigma_x,\sigma_y,\sigma_z$ 表示绕三个坐标轴的半周转动，分别为
\begin{equation}
\operatorname{diag}\left(1,-1,-1\right),\qquad\operatorname{diag}\left(-1,1,-1\right),\qquad\operatorname{diag}\left(-1,-1,1\right).
\end{equation}
此处这些符号不是 Pauli 矩阵。$\mathrm{D}_3=\left\{e,d,f,a,b,c\right\}$ 的类和表示已在前文讨论。

原稿中的直积分解 $\mathrm{D}_{4n+2}=\mathrm{D}_{2n+1}\otimes\mathrm{C}_2$ 可如下验证。
令 $m=2n+1$ 为奇数，取 $\mathrm{D}_{2m}$ 的生成元 $r_0,s_0$，满足
$r_0^{2m}=s_0^2=e$、$s_0r_0s_0=r_0^{-1}$。
子群 $H=\langle r_0^2,s_0\rangle$ 同构于 $\mathrm{D}_m$，
而 $F=\left\{e,r_0^m\right\}$ 同构于 $\mathrm{C}_2$。
由于 $s_0r_0^ms_0=r_0^{-m}=r_0^m$，$F$ 与所有群元对易；
$m$ 为奇数，又保证 $r_0^m$ 不在 $H$ 中，所以 $H\cap F=\left\{e\right\}$。
两子群的乘积有 $\left(2m\right)\cdot2=4m$ 个元素，等于全群阶，因此
$\mathrm{D}_{2m}=H\otimes F$，给出所述分解。
\end{example}

\section{群函数、群代数与正则表示}
\subsection{群函数和群代数}
有限群 $G$ 上的复值函数 $F:G\to\mathbb{C}$ 称为群函数，其内积取为
\begin{equation}
\left\langle F_1,F_2\right\rangle=\frac{1}{\left|G\right|}\sum_{g\in G}F_1\left(g\right)^*F_2\left(g\right).
\end{equation}
在形式线性组合
\begin{equation}
x=\sum_{g\in G}x_g g
\end{equation}
上把群乘法双线性延拓，得到复群代数 $\mathbb{C}\left[G\right]$。它既是 $\left|G\right|$ 维线性空间，也是结合代数。

\subsection{左正则表示和内禀正则表示}
在群代数基 $\Ket{g}$ 上定义
\begin{equation}
L_s\Ket{g}=\Ket{sg},
\end{equation}
得到左正则表示，其矩阵元和特征标为
\begin{align}
\left[L_s\right]_{hg}&=\delta_{h,sg},\\
\chi_{\mathrm{reg}}\left(s\right)&=\begin{cases}\left|G\right|,&s=E,\\0,&s\ne E.
\end{cases}
\end{align}
右乘满足相反的乘法次序。定义
\begin{equation}
\widetilde R_s\Ket{g}=\Ket{gs^{-1}},
\end{equation}
则 $\widetilde R_s\widetilde R_t=\widetilde R_{st}$，称为内禀正则表示，并且
\begin{equation}
\left[L_s,\widetilde R_t\right]=0.
\end{equation}

\section{新表示的构成}
\subsection{商群、直乘群与表示直乘}
若 $N\trianglelefteq G$，$G/N$ 的表示 $D$ 可提升为
\begin{equation}
\widetilde D\left(g\right)=D\left(gN\right).
\end{equation}
反过来，$G$ 的表示能够下降到 $G/N$ 的充要条件是 $N$ 中每个元素都表示为单位矩阵。

若 $D^{\left(i\right)},D^{\left(j\right)}$ 是 $G$ 的表示，则
\begin{equation}
D\left(g\right)=D^{\left(i\right)}\left(g\right)\otimes D^{\left(j\right)}\left(g\right)
\end{equation}
仍为表示，且
\begin{equation}
\chi\left(g\right)=\chi^{\left(i\right)}\left(g\right)\chi^{\left(j\right)}\left(g\right).
\end{equation}
不可约表示 $D^{\left(k\right)}$ 的重数为
\begin{equation}
a_k=\frac{1}{\left|G\right|}\sum_{g\in G}\chi^{\left(k\right)}\left(g\right)^*\chi^{\left(i\right)}\left(g\right)\chi^{\left(j\right)}\left(g\right).
\end{equation}
若 $G=H_1\otimes H_2$，则全部不可约表示为
\begin{equation}
D^{\left(i,j\right)}\left(h_1,h_2\right)=D^{\left(i\right)}_1\left(h_1\right)\otimes D^{\left(j\right)}_2\left(h_2\right).
\end{equation}

\subsection{限制表示、诱导表示与 Frobenius 互反律}
设 $H\leq G$。把 $G$ 的表示限制在 $H$ 上，得到 $\operatorname{Res}^G_H D$。反过来，取左陪集代表元 $t_1,\ldots,t_q$，其中 $q=\left[G:H\right]$。若 $d$ 是 $H$ 的表示，且
\begin{equation}
g t_j=t_i h_{ij}\left(g\right),qquad h_{ij}\left(g\right)\in H,
\end{equation}
则诱导表示定义为
\begin{equation}
\left[\operatorname{Ind}_H^G d\left(g\right)\right]_{ij}=
\begin{cases}d\left(h_{ij}\left(g\right)\right),&g t_jH=t_iH,\\0,&\text{其他情形}.
\end{cases}
\end{equation}
它的维数为 $q\dim d$。Frobenius 互反律为
\begin{equation}
\left\langle \operatorname{Res}^G_H\chi_D,\chi_d\right\rangle_H
=\left\langle \chi_D,\operatorname{Ind}_H^G\chi_d\right\rangle_G.
\end{equation}

\section{投影算符与群代数分解}
设 $D^{\left(i\right)}$ 是 $m_i$ 维不可约幺正表示，$P_g$ 是空间上的群作用。矩阵元投影算符为
\begin{equation}
P^{\left(i\right)}_{\mu\nu}=\frac{m_i}{\left|G\right|}\sum_{g\in G}D^{\left(i\right)}_{\mu\nu}\left(g\right)^*P_g,
\end{equation}
并满足
\begin{equation}
P^{\left(i\right)}_{\mu\nu}P^{\left(j\right)}_{\rho\sigma}
=\delta_{ij}\delta_{\nu\rho}P^{\left(i\right)}_{\mu\sigma}.
\end{equation}
特征标投影算符为
\begin{equation}
P^{\left(i\right)}=\frac{m_i}{\left|G\right|}\sum_{g\in G}\chi^{\left(i\right)}\left(g\right)^*P_g,
\end{equation}
它们满足
\begin{align}
P^{\left(i\right)}P^{\left(j\right)}&=\delta_{ij}P^{\left(i\right)},\\
\sum_iP^{\left(i\right)}&=E.
\end{align}
在复数域上，有限群群代数分解为
\begin{equation}
\mathbb{C}\left[G\right]\cong\bigoplus_i M_{m_i}\left(\mathbb{C}\right),
\qquad \left|G\right|=\sum_i m_i^2.
\end{equation}
中心本原幂等元为
\begin{equation}
e_i=\frac{m_i}{\left|G\right|}\sum_{g\in G}\chi^{\left(i\right)}\left(g^{-1}\right)g,
\end{equation}
且 $e_ie_j=\delta_{ij}e_i$、$\sum_i e_i=E$。因此左正则表示中，第 $i$ 个不可约表示出现 $m_i$ 次。


